\section{RESULTADOS}

%Debe presentar los resultados del proceso de control de calidad, según lo definido en la meto-dología. Por ejemplo, en un proyecto de desarrollo de software, esta sección debe incluir los resultados de pruebas estándar de software (unidad, funcional, sistema, aceptación, etc.). Es importante que sean coherentes con el objetivo del proyecto y la metodología utilizada para su desarrollo.
%Este capítulo debe incluir un análisis de los resultados obtenidos y las conclusiones derivadas de dicho análisis.

\subsection{Límites de la evidencia}
El periodo reservado se observó durante el desarrollo del sistema. El protocolo, la partición de la reserva, los umbrales y las semillas se fijaron antes de la corrida final y no se ajustaron con sus resultados. Los productos que no se usaron durante el desarrollo se reportan también por separado.

La regla que exige al menos dos familias evaluables deja fuera de la competencia a los productos en que solo una familia produce un modelo válido. Por eso, los resultados describen la población de productos en que al menos dos familias producen un modelo válido. Para cada clase se reportan los productos excluidos con su causa y se compara su perfil de demanda previo a la reserva con el de los productos conservados.

El bloque de prueba no contiene la temporada de fin de año. El desempeño de diciembre se describe con el bloque de selección, sin pruebas estadísticas.
